\section{Jump loci and agreement of the segment enumeration}\label{si:loci}

\paragraph{Agreement.} The declared hold class couples the stages of the blind window, so its value is not computed by the stagewise recursion over $A=\{\pm1\}$---that recursion computes the sequential value, which equals the open-loop value \meqref{eq:ol} and is one on the timing cells. The segment form of Proposition~\mref{prop:alpha} enumerates the two admissible holds and evaluates the surviving mass of each; Table~\ref{tab:grid} records that this enumeration returns \meqref{eq:hold} on all $192$ combinations. The class constraint is load-bearing: Corollary~\mref{cor:timing} is a statement about which class the timing certificate speaks for.

\paragraph{Jump loci.} Fix $\Tobs$, $k\ge1$ and $w$, and read \meqref{eq:hold} as a function of the real parameter $z_0$. It is piecewise constant with exactly two jumps: at $z_0=1$, from $0$ to $\max(w,1-w)$, and at $z_0=1+\ell$, $\ell=\min(k,\Tobs)$, from $\max(w,1-w)$ to $1$; at the symmetric prior both jumps have size $\tfr12$. For $k\le\Tobs$ the second locus is $z_0=1+k$---an off-grid consequence of the closed form (for $k=2$ and $\Tobs\ge2$ the jump is at $z_0=3$, not at $2$), verified on the extended grid. The witness family changes at the same loci: at the one-step horizon both mirror witnesses $(1,0)$ and $(0,1)$ attain the value $\max(w,1-w)$ below $z_0=2$ and the merged witness $(1,1)$ attains $1$ at and above it; at horizon $k$ the change sits at $z_0=1+\ell$. These are jumps in a parameter of the instance; as a function of the belief $b$ the value is piecewise linear and convex (Proposition~\mref{prop:alpha}), its kinks being the usual ones of the $\alpha$-vector representation.

\paragraph{Deficit profile.} On the $42$ nonviable cells the hold-class deficit at horizon $k$ is $\tfr12$ if $k>z_0-1$ and $0$ otherwise (for $k\le\Tobs$), and $\tfr12$ for all $k>\Tobs$: on the timing cells $0$ at $k=1$ and $\tfr12$ at $k\ge2$, on the common-action cells $\tfr12$ at every $k\ge1$. The deficit is non-decreasing in $k$ (Proposition~\mref{prop:deficit}(iii)) and frozen from $k=\Tobs$ on (Proposition~\mref{prop:freeze}); the open-loop deficit is frozen from $k=1$. Both freezes are verified by enumeration to horizon $40$.
